\section{Conclusion}
This study provides fairly comprehensive empirical evidence on the inherent vulnerabilities of Retrieval-Augmented Generation systems when exposed to semantic mutations. By evaluating across different model sizes and prompt configurations, we identified what we call the Inverse Scaling Law of Semantic Poisoning --- showing that, contrary to common intuition, larger models are actually more vulnerable when they lack proper safeguards. We also demonstrated the high reliability and consistency of the LLM-as-a-Judge approach when applied across both public and private datasets.

Finally, our results show that semantic attacks tend to trigger misinformation acceptance rather than hallucination, and that Prompt Engineering alone is not enough to secure RAG deployments in practice. In summary, while prompt design can help reduce vulnerabilities to some extent, architectural improvements remain essential to ensuring the integrity of LLM applications against adversarial data streams.
